% Fragment partage : inclus par socle-compte.tex (dossier autonome)
% et par dossier-conception-alanya.tex (dossier client assemble).
% Ne pas y mettre de preambule ni de \part : c'est l'appelant qui les pose.

\chapter{Ce que l'on ajoute, et pourquoi}

\section{La situation aujourd'hui}

Dans Alanya, tous les comptes sont identiques. Un commerçant, une personnalité publique et
un utilisateur ordinaire ont exactement le même profil et la même apparence dans une liste
de conversations. Rien ne permet de distinguer l'entreprise officielle d'un imitateur qui
aurait pris le même nom et la même photo.

C'est ce manque que les trois fonctionnalités demandées viennent combler --- et c'est
pourquoi elles partagent la même fondation.

\section{Deux questions, pas une}

L'erreur naturelle serait de créer une liste de « genres de comptes » : personnel,
business, vérifié, officiel, business vérifié\ldots{} Cette liste devient ingérable dès
qu'on ajoute un cas.

On sépare donc deux questions \textbf{indépendantes} :

\begin{center}
\begin{tabularx}{\linewidth}{@{}p{4.2cm}X@{}}
\toprule
\thd{Quel genre de compte ?} & Personnel, business, ou officiel. C'est ce que le compte
\emph{est}. \\
\thd{A-t-il été vérifié ?} & Aucune démarche, en attente, vérifié, refusé, révoqué, ou
expiré. C'est ce qu'Alanya \emph{atteste} à son sujet. \\
\bottomrule
\end{tabularx}
\end{center}

Ces deux questions ne dépendent pas l'une de l'autre. Un commerce peut exister sans être
vérifié. Une personne privée peut être vérifiée sans être un commerce --- c'est le cas
d'une personnalité publique. Le compte officiel Alanya, lui, est la combinaison
« officiel + vérifié d'office ».

\begin{note}[Le bénéfice concret]
Poser deux questions séparées plutôt qu'une seule liste, c'est ce qui permet d'ajouter
demain un nouveau genre de compte, ou un nouvel état de vérification, sans toucher à
l'autre axe ni refaire les écrans.
\end{note}

\section{Ce que couvre l'étape 1 --- et ce qu'elle ne couvre pas}

\begin{tabularx}{\linewidth}{@{}p{5.4cm}X@{}}
\toprule
\thd{Dans cette étape} & Savoir \emph{qui} est un compte : son genre, son état de
vérification, et le badge qui en découle partout dans l'application. \\
\thd{Étapes suivantes} & Ce que ce compte \emph{contient} : la fiche business (horaires,
adresse, catalogue, liens), le dossier de certification (pièces justificatives,
approbation, abonnement), et la diffusion du compte officiel. \\
\bottomrule
\end{tabularx}

La raison de cette coupure est simple : le genre et l'état d'un compte suffisent déjà à
afficher le bon badge partout. Le contenu détaillé d'une fiche business dépend de règles
qu'on n'a pas encore écrites --- le figer maintenant obligerait à tout refaire à l'étape
suivante.

\chapter{Ce que l'utilisateur voit}

\section{Le principe du badge}

Un badge qui combinerait « c'est un commerce » et « c'est vérifié » en un seul dessin
devient illisible à la taille où il apparaît réellement --- environ 12 pixels, à côté d'un
nom dans une liste de conversations. Et deux badges côte à côte mangent la largeur du nom,
qui est déjà tronqué sur un petit écran.

La solution retenue sépare les deux informations sur deux canaux visuels distincts :

\begin{center}
\begin{tabularx}{0.92\linewidth}{@{}p{4.6cm}X@{}}
\toprule
\thd{Le glyphe dit \emph{quoi}} & Une coche pour une personne notable, un panier pour un
commerce, la marque Alanya pour le compte officiel. \\
\thd{Le contenant dit \emph{vérifié ou pas}} & Un glyphe nu = déclaré par le compte
lui-même. Un glyphe posé dans un \textbf{sceau festonné} \badgeSceauNu[0.9] = attesté par
Alanya. \\
\bottomrule
\end{tabularx}
\end{center}

\subsection*{Pourquoi un sceau festonné, et pas un rond}

Un disque plein ne se distingue d'aucune autre icône circulaire de l'interface : avatars,
puces de comptage, indicateurs de présence. À 12~pixels, un badge rond devient un point
coloré parmi d'autres points colorés.

Le feston à douze lobes, lui, a une silhouette reconnaissable même quand le glyphe qu'il
contient n'est plus lisible. C'est la propriété qui compte : l'utilisateur qui distingue
la \emph{forme} sait qu'Alanya s'est prononcé sur ce compte, sans avoir à lire le
détail. Et cette forme est bien plus difficile à imiter avec un caractère Unicode dans un
nom d'affichage qu'un simple rond ou une simple coche.

\section{Les cinq rendus}

\begin{center}
\setlength{\extrarowheight}{7pt}
\begin{tabular}{@{}llc@{\hspace{1.6cm}}c@{}}
\toprule
\thd{Genre} & \thd{État} & \thd{À taille réelle} & \thd{Agrandi} \\
\midrule
Personnel & non vérifié & --- & --- \\
Personnel & vérifié     & \badgePastilleCoche[1]  & \badgePastilleCoche[2.2] \\
Business  & déclaré     & \badgePanier[1]         & \badgePanier[2.2] \\
Business  & vérifié     & \badgePastillePanier[1] & \badgePastillePanier[2.2] \\
Officiel  & ---         & \badgeMarque[1]         & \badgeMarque[2.2] \\
\bottomrule
\end{tabular}
\end{center}

\medskip

La colonne « à taille réelle » est là pour être jugée telle quelle : c'est exactement la
taille à laquelle le badge apparaîtra à côté d'un nom. Si un rendu n'est pas
distinguable à cette échelle, il doit être redessiné avant l'implémentation.

\begin{note}[Pourquoi une couleur différente pour l'officiel]
Le sceau du compte officiel est doré, pas indigo. Il n'existe qu'un seul compte
officiel dans toute l'application : lui donner une couleur à part évite qu'il soit confondu
avec un compte simplement vérifié, et rend l'usurpation visuellement plus difficile.
\end{note}

\section{Où le badge apparaît}

Le même badge, dessiné par le même composant, dans quatre contextes :

\begin{itemize}[leftmargin=1.4em]
  \item la \textbf{liste de conversations}, juste après le nom ;
  \item l'\textbf{en-tête d'une conversation} ouverte ;
  \item la \textbf{page de profil} du compte ;
  \item les \textbf{résultats de recherche} et les listes de contacts.
\end{itemize}

\begin{center}
  \imageOuCadre{maquettes/liste-conversations-badges.png}{0.86\linewidth}%
    {Liste de conversations montrant les cinq variantes de badge en situation\\
     \texttt{maquettes/liste-conversations-badges.png}}
\end{center}

\begin{note}[Maquette interactive]
La version vivante de cette maquette permet de faire varier la taille du badge de 10 à
28~pixels pour juger la lisibilité, et montre la conversation officielle avec son bandeau
de lecture seule :\\
\href{https://claude.ai/code/artifact/029627c6-5dae-4b2a-a86b-3e28c899d824}%
{\texttt{claude.ai/code/artifact/029627c6}}\\
Source versionnée : \code{docs/architecture/maquettes/liste-conversations-badges.html}
\end{note}

\section{Pourquoi la distinction déclaré / vérifié est critique}

Se déclarer « commerce » prend trente secondes et n'engage personne : le compte l'affirme
lui-même, personne ne l'a contrôlé. Être vérifié suppose au contraire un dossier examiné
par Alanya.

Si les deux badges avaient le même poids visuel, n'importe qui créerait un compte business
et hériterait d'une caution qui ne lui a jamais été accordée. C'est exactement le scénario
qu'un arnaqueur cherche.

\begin{alerte}[La règle qui en découle]
Le sceau festonné est la signature d'Alanya. Rien d'autre dans l'interface ne doit pouvoir lui
ressembler --- ni une icône décorative, ni un emoji dans un nom d'utilisateur.
\end{alerte}

\section{Protéger le badge des contournements}

Un système de badge peut être contourné sans jamais être attaqué techniquement : il suffit
qu'un utilisateur s'appelle « Alanya~\checkmark{} » pour que le nom affiché imite la
certification.

Deux protections sont donc posées dès cette étape :

\begin{enumerate}[leftmargin=1.6em]
  \item les caractères de type coche, étoile et sceau sont \textbf{refusés} dans le nom et
        le pseudo ;
  \item les variantes du nom « Alanya » sont \textbf{réservées} et ne peuvent pas être
        prises par un compte ordinaire.
\end{enumerate}

\chapter{Les règles de vie d'un compte}

\section{Devenir un compte business, et revenir en arrière}

Le passage d'un compte personnel à un compte business est \textbf{réversible}. Les règles
proposées :

\begin{center}
\begin{tabularx}{\linewidth}{@{}p{5.2cm}X@{}}
\toprule
\thd{Personnel $\rightarrow$ business} & Libre, à l'initiative de l'utilisateur. Devenir
commerce est déclaratif : cela n'atteste de rien et ne demande donc aucune validation. \\
\thd{Business $\rightarrow$ personnel} & Libre également. La fiche business est
\textbf{conservée mais masquée}, pas supprimée : un retour en arrière ne doit pas détruire
un catalogue de produits saisi pendant des semaines. \\
\thd{Une vérification en cours} & La demande est annulée : elle portait sur une identité
d'entreprise qui n'est plus revendiquée. \\
\thd{Une vérification obtenue} & Elle est \textbf{révoquée}, pas suspendue. Le badge
atteste d'une entreprise ; si le compte n'en est plus une, l'attestation n'a plus d'objet.
Un nouveau dossier sera nécessaire pour la retrouver. \\
\bottomrule
\end{tabularx}
\end{center}

\begin{note}[À confirmer avec le client]
La bascule business $\rightarrow$ personnel pourrait aussi être soumise à validation d'un
administrateur, pour éviter qu'un commerce échappe à un signalement en repassant en
personnel. La règle proposée ci-dessus est la plus simple ; l'alternative est plus sûre.
\end{note}

\section{La vie d'une vérification}

\begin{center}
\begin{tikzpicture}[
  x=1cm, y=1cm,
  etat/.style={draw=AlanyaLine, fill=AlanyaSurface, rounded corners=3pt,
               font=\sffamily\small, inner sep=6pt, minimum height=0.8cm,
               text=AlanyaInk},
  etatOk/.style={etat, draw=AlanyaSuccess, line width=0.9pt, fill=AlanyaSuccess!8},
  etatKo/.style={etat, draw=danger, fill=danger!6},
  fleche/.style={-{Stealth[length=2.2mm]}, draw=AlanyaGray, semithick},
  etiq/.style={font=\sffamily\scriptsize, text=AlanyaGray,
               fill=white, inner sep=1.5pt}
]
  \node[etat]   (aucun)   at (0,0)      {Aucune démarche};
  \node[etat]   (attente) at (4.3,0)    {En attente};
  \node[etatOk] (ok)      at (8.3,0)    {Vérifié};
  \node[etatKo] (revoque) at (11.8,0)   {Révoqué};
  \node[etatKo] (refuse)  at (1.7,-2.1) {Refusé};
  \node[etatKo] (expire)  at (8.3,-2.1) {Expiré};

  % Chemin nominal
  \draw[fleche] (aucun)   -- node[etiq, above] {dépôt}       (attente);
  \draw[fleche] (attente) -- node[etiq, above] {approbation} (ok);
  \draw[fleche] (ok)      -- node[etiq, above] {abus}        (revoque);
  \draw[fleche] (ok)      -- node[etiq, right] {échéance}    (expire);

  % Rejet
  \draw[fleche] (attente.south west) -- node[etiq, pos=0.55, above left, inner sep=1pt]
    {rejet} (refuse.north east);

  % Retours : deux couloirs qui ne se croisent pas
  \draw[fleche] (refuse.north west) -- node[etiq, pos=0.5, above left, inner sep=1pt]
    {nouveau dossier} (aucun.south);

  \draw[fleche] (expire.west) -- (4.3,-2.1)
    node[etiq, pos=0.55, below] {renouvellement} -- (attente.south);
\end{tikzpicture}
\end{center}

\section{L'expiration et sa relance}

Une vérification n'est pas éternelle : elle a une date de fin. Le comportement attendu :

\begin{center}
\begin{tabularx}{\linewidth}{@{}p{3.4cm}X@{}}
\toprule
\thd{À J\,--\,30} & Le titulaire est prévenu qu'il doit renouveler, par trois canaux à la
fois : une notification sur son téléphone, un e-mail, et un avertissement visible dans
l'application. L'avertissement n'est envoyé qu'une fois. \\
\thd{À l'échéance} & La vérification est \textbf{suspendue}. Le badge disparaît. Le compte
et son contenu restent intacts --- seule l'attestation tombe. \\
\thd{Après} & Le titulaire peut déposer un nouveau dossier pour retrouver son badge. \\
\bottomrule
\end{tabularx}
\end{center}

\begin{note}[Suspendu, pas supprimé]
Une vérification expirée est un état distinct d'un refus. Le compte a bien été vérifié par
le passé ; il a simplement laissé passer la date. Cette nuance compte pour le service
client, et pour ne pas retraiter un dossier déjà instruit.
\end{note}

\section{Le compte officiel}

Le compte officiel est le compte au nom de l'application. Ses règles :

\begin{itemize}[leftmargin=1.4em]
  \item il est \textbf{créé par un administrateur}, jamais par une inscription ordinaire ;
  \item sa photo de profil est le \textbf{logo Alanya} ;
  \item il porte le badge officiel, distinct du badge de vérification ;
  \item les utilisateurs \textbf{ne peuvent pas lui répondre}. Ses messages sont des
        diffusions, pas des conversations.
\end{itemize}

Concrètement, dans une conversation avec le compte officiel, la zone de saisie est
remplacée par un bandeau explicatif, et les actions « répondre » et « réagir » n'apparaissent
pas sur ses messages.

\begin{alerte}[Ce n'est pas qu'une question d'affichage]
Masquer la zone de saisie ne suffit pas : un client modifié pourrait envoyer le message
quand même. Le refus doit être appliqué par le serveur. La partie II décrit où.
\end{alerte}

\section{Un compte business ne peut pas être administrateur}

Décision prise : les rôles d'administration de la plateforme et les genres de compte sont
mutuellement exclusifs. Un compte business ou officiel ne peut pas recevoir de droits
d'administration, et réciproquement.

La raison est de sécurité : un compte visible du public, susceptible d'être usurpé ou
compromis, ne doit pas ouvrir la porte au back-office.
